% =====================================================================
%  07_trabajo-futuro.tex — Pendientes y siguientes prototipos
% =====================================================================
\chapter{Trabajo a futuro}
\label{cap:futuro}

Los resultados del Prototipo~1 definen el trabajo de los siguientes
prototipos. Se agrupa por prototipo según el plan del
capítulo~\ref{cap:metodologia}.

\section{Prototipo 2: operación continua}

\begin{itemize}
  \item \textbf{Instalación en el gabinete} del patio con la Jetson de
        16~GB y fuente propia~\cite{rep-tecnico}.
  \item \textbf{Arranque automático}: hoy los servicios tienen
        \texttt{restart: "no"} y el sistema no arranca solo al encender la
        Jetson~\cite{rep-tecnico}. Debe probarse un reinicio controlado con
        alguien en sitio.
  \item \textbf{Entrega de avisos verificada}: dar seguimiento al
        desbloqueo de Telegram en la red del patio~\cite{rep-envivo},
        confirmar la solución del incidente de WhatsApp~\cite{rep-whatsapp}
        y avisar cuando un canal deje de entregar mensajes.
  \item \textbf{Monitoreo}: avisar cuando la cámara o un servicio dejen de
        funcionar.
\end{itemize}

\section{Prototipo 3: mejora de la lectura}

\begin{itemize}
  \item \textbf{Medir la exactitud} de cada etapa (detección, zona y
        dígitos) con una muestra de cruces revisada contra la grabación.
  \item \textbf{Ampliar los datos}: etiquetar los recortes guardados en el
        patio y versionar el conjunto de datos.
  \item \textbf{Reentrenar M2} con más ejemplos de vistas delantera y
        trasera, donde hoy casi no encuentra la zona del número.
  \item \textbf{Exportar} los nuevos modelos a TensorRT en la propia Jetson
        y comparar contra la versión anterior.
\end{itemize}

\section{Pendientes del proyecto}

\begin{itemize}
  \item \textbf{Una sola fuente de verdad del código}: unificar la versión
        del repositorio, el paquete \texttt{SARU\_V4} y la copia de la
        Jetson, y trabajar con ramas y revisión de cambios.
  \item \textbf{Validar con video real} la línea virtual de varios puntos
        (LVC-13).
  \item \textbf{Agregar el simulador al repositorio}, para que las pruebas
        LVC puedan repetirse en cada cambio.
  \item \textbf{Corregir la configuración}: rutas de datos por defecto que
        no coinciden con la Jetson, variables de Telegram exigidas aunque
        los avisos estén apagados y documentación del repositorio que cita
        scripts inexistentes~\cite{rep-tecnico}.
  \item \textbf{Validar las historias de usuario} con el personal del patio
        y cuantificar la problemática.
\end{itemize}
